\labsection{7}{Support vector machines}{sec:lab07-svm}

\begin{labproject}{Tune a margin-based classifier inside a leakage-safe pipeline}
\textbf{Coverage.} Chapter 14.\\
\textbf{Core data.} Breast Cancer Wisconsin.\\
\textbf{Goal.} Standardize features inside the pipeline, tune SVM capacity with cross-validation, and interpret how $C$ and kernel settings change flexibility.

\begin{labmode}
This is the final supervised model-family lab. Reuse the validation discipline established in Lab 3 rather than inventing a new evaluation process for SVMs.
\end{labmode}

\smallskip\noindent\textit{Implementation note.} Code is shown in notebook-sized blocks rather than as one complete listing. Run the blocks in order; each framed block corresponds to a canonical Jupyter code cell.\par

\medskip\noindent\textbf{Part A --- Standardized SVM tuning.}\par
\lstinputlisting[style=mlpython]{jupyter_cells/lab07-support_vector_machines/block01.py}
\lstinputlisting[style=mlpython]{jupyter_cells/lab07-support_vector_machines/block02.py}
\lstinputlisting[style=mlpython]{jupyter_cells/lab07-support_vector_machines/block03.py}
\lstinputlisting[style=mlpython]{jupyter_cells/lab07-support_vector_machines/block04.py}
\end{labproject}

\begin{labdeliverable}
Submit the cross-validated SVM settings, validation evidence for the selected configuration, and a short explanation of why scaling and hyperparameter selection must remain inside the validation procedure.
\end{labdeliverable}
